\id{ҒТАМР 65.09.05}{}

\begin{header}
\swa{Д.Б.~Курмангалиева, А.А.~Дарибаева, Г.Т.~Юсупова, Ж.К.~Шадьярова, Г.Г.~Жуматова}{БАЛЫҚ ӨНІМДЕРІНІҢ САПАСЫ МЕН ҚАУІПСІЗДІГІН НОРМАТИВТІК РЕТТЕУ ЖӘНЕ ХАЛЫҚАРАЛЫҚ СТАНДАРТТАРМЕН ҮЙЛЕСТІРУ}

\tsp{1}Д.Б.~Курмангалиева,
\tsp{1}А.А.~Дарибаева,
\tsp{1}Г.Т.~Юсупова\envelope,
\tsp{1}Ж.К.~Шадьярова,
\tsp{2}Г.Г.~Жуматова
\end{header}

\begin{affil}
\tsp{1}С.Сейфуллин атындағы Қазақ агротехникалық зерттеу университеті КеАҚ, Астана, Қазақстан,

\tsp{2}Х. Досмұхамедов атындағы Атырау университеті КеАҚ, Атырау, Қазақстан,

\corrauthor{Корреспондент-автор: gauhar\_20\_90@mail.ru}
\end{affil}

Халықаралық деңгейде азық-түлік қауіпсіздігіне қойылатын талаптардың
күшеюі жағдайында ұлттық нормативтік реттеу жүйелерін халықаралық
стандарттармен үйлестіру өзекті мәселеге айналуда. Балық өнімдері,
әсіресе ready-to-eat (RTE) санатына жататын өнімдер микробиологиялық
қауіпсіздікке, суық тізбекті бақылауға және шикізаттың қадағалануына
қатысты жоғары талаптармен сипатталады.

Зерттеудің мақсаты -- Еуразиялық экономикалық одақ елдеріндегі балық
өнімдерінің қауіпсіздігін нормативтік реттеу жүйесінің үйлесімділік
деңгейін Еуропалық одақ талаптарымен және Codex Alimentarius халықаралық
стандарттарымен салыстыра отырып бағалау.

Зерттеу барысында нормативтік құжаттарды салыстырмалы-құқықтық талдау,
нормативтік талаптардың контент-талдауы, реттеуші блоктарды сараптамалық
ранжирлеу, gap-талдау және нормативтік жүйенің үйлесімділігін интегралды
бағалау әдістерін қамтитын ғылыми әдістер кешені қолданылды. Зерттеу
шеңберінде ұлттық нормативтік жүйенің халықаралық талаптарға сәйкестік
деңгейін сандық тұрғыдан бағалауға мүмкіндік беретін IH-FISH (Index of
Harmonization for Fish Safety Regulation) интегралды үйлесімділік
индексінің моделі әзірленді.

Зерттеу нәтижелері нормативтік талаптардың ең жоғары үйлесімділік
деңгейі микробиологиялық қауіпсіздік және өндірістің санитарлық
талаптары саласында байқалатынын көрсетті. Сонымен қатар суық тізбекті
цифрлық мониторингтеу, өнімнің қадағалану жүйелері және антифрод-бақылау
механизмдері саласында елеулі айырмашылықтар анықталды. Интегралды
үйлесімділік индексінің есептеулері ЕАЭО нормативтік жүйесінің
халықаралық талаптарға сәйкестігі орташа деңгейде екенін көрсетті.

Алынған нәтижелер балық өнімдерінің қауіпсіздігін нормативтік реттеу
жүйесін жетілдіруде, жеткізу тізбектерін цифрлық мониторингтеу
механизмдерін дамытуда және өңделген балық өнімдерінің халықаралық
нарықтағы экспорттық бәсекеге қабілеттілігін арттыруда қолданылуы
мүмкін.

{\bfseries Түйін сөздер:} балық өнімдері, тағам қауіпсіздігі, стандарттарды
үйлестіру, антифрод-бақылау, gap-талдау, HACCP.

\begin{header}
НОРМАТИВНОЕ РЕГУЛИРОВАНИЕ КАЧЕСТВА И БЕЗОПАСНОСТИ РЫБНОЙ ПРОДУКЦИИ И ГАРМОНИЗАЦИЯ С МЕЖДУНАРОДНЫМИ СТАНДАРТАМИ

\tsp{1}Д.Б.~Курмангалиева,
\tsp{1}А.А.~Дарибаева,
\tsp{1}Г.Т.~Юсупова\envelope,
\tsp{1}Ж.К.~Шадьярова,
\tsp{2}Г.Г.~Жуматова
\end{header}

\begin{affil}
\tsp{1}НАО «Казахский агротехнический исследовательский университет им. С.Сейфуллина», Астана, Казахстан,

\tsp{2}НАО «Атырауский университет им. Х. Досмухамедова», Атырау, Казахстан,

e-mail: gauhar\_20\_90@mail.ru
\end{affil}

\emph{В условиях} усиления международных требований к безопасности
пищевых продуктов особую актуальность приобретает гармонизация
национальных систем нормативного регулирования с международными
стандартами. Рыбная продукция, особенно изделия категории ready-to-eat
(RTE), характеризуется повышенными требованиями к микробиологической
безопасности, контролю холодовой цепи и прослеживаемости сырья.

Целью настоящего исследования является оценка уровня гармонизации
нормативного регулирования безопасности рыбной продукции в странах
Евразийского экономического союза с требованиями Европейского союза и
международными стандартами Codex Alimentarius.

В работе использован комплекс научных методов, включающий
сравнительно-правовой анализ нормативных документов, контент-анализ
регуляторных требований, экспертное ранжирование регуляторных блоков,
gap-анализ и метод интегральной оценки гармонизации нормативной системы.
В рамках исследования была разработана модель интегрального индекса
гармонизации IH-FISH (Index of Harmonization for Fish Safety
Regulation), позволяющая количественно оценить степень соответствия
национальной нормативной системы международным требованиям.

Результаты исследования показали, что наиболее высокий уровень
гармонизации нормативных требований наблюдается в области
микробиологической безопасности и санитарных требований производства.
Существенные различия выявлены в сфере цифрового мониторинга холодовой
цепи, систем прослеживаемости продукции и механизмов антифрод-контроля.
Расчёт интегрального индекса гармонизации показал средний уровень
соответствия нормативной системы ЕАЭС международным требованиям.

Полученные результаты могут быть использованы для совершенствования
системы нормативного регулирования безопасности рыбной продукции,
развития механизмов цифрового мониторинга цепей поставок и повышения
экспортной конкурентоспособности переработанных рыбных продуктов на
международных рынках.

{\bfseries Ключевые слова:} рыбная продукция, безопасность пищевых
продуктов, гармонизация стандартов, антифрод-контроль, gap-анализ,
HACCP.

\begin{header}
{\bfseries REGULATORY FRAMEWORK FOR QUALITY AND SAFETY OF FISH PRODUCTS AND HARMONIZATION WITH INTERNATIONAL STANDARDS}

\tsp{1}D.~Kurmangaliyeva,
\tsp{1}A.~Daribayeva,
\tsp{1}G.~Yussupova\envelope,
\tsp{1}Zh.~Shadyarova,
\tsp{2}G.~Zhumatova
\end{header}

\begin{affil}
\tsp{1}NP JSC «S. Seifullin Kazakh Agrotechnical Research University», Astana, Kazakhstan

\tsp{2}NP JSC «Atyrau University named after Khalel Dosmukhamedov» Atyrau, Kazakhstan

e-mail: gauhar\_20\_90@mail.ru
\end{affil}

In the context of increasing international requirements for food safety,
the harmonization of national regulatory systems with international
standards has become particularly relevant. Fish products, especially
ready-to-eat (RTE) products, are characterized by increased requirements
for microbiological safety, cold chain control, and traceability of raw
materials.

The aim of this study is to assess the level of harmonization of the
regulatory framework for fish product safety in the countries of the
Eurasian Economic Union with the requirements of the European Union and
the international standards of Codex Alimentarius.

The study employed a комплекс of scientific methods, including
comparative legal analysis of regulatory documents, content analysis of
regulatory requirements, expert ranking of regulatory blocks, gap
analysis, and the method of integral assessment of regulatory
harmonization. Within the framework of the study, a model of the IH-FISH
(Index of Harmonization for Fish Safety Regulation) integral
harmonization index was developed, allowing a quantitative assessment of
the compliance of national regulatory systems with international
requirements.

The results of the study showed that the highest level of harmonization
of regulatory requirements is observed in the field of microbiological
safety and sanitary requirements of production. Significant differences
were identified in the areas of digital monitoring of the cold chain,
product traceability systems, and anti-fraud control mechanisms. The
calculation of the integral harmonization index demonstrated a moderate
level of compliance of the EAEU regulatory system with international
requirements.

The obtained results can be used to improve the regulatory system for
fish product safety, develop digital monitoring mechanisms for supply
chains, and enhance the export competitiveness of processed fish
products in international markets.

{\bfseries Keywords:} fish products; food safety, harmonization of
standards, anti-fraud control, gap analysis, HACCP.

\begin{multicols}{2}
{\bfseries Кіріспе.} Соңғы онжылдықтарда әлемдік балық өнімдері өндірісі
тұрақты өсім көрсетіп келеді, бұл аквакультураның белсенді дамуы, өңдеу
технологияларының жетілдірілуі және жоғары тағамдық әрі биологиялық
құндылығы бар өнімдерге тұтынушылық сұраныстың артуымен байланысты.
Балық өнімдері әлемдік азық-түлік нарығы құрылымында маңызды орын алады
және толыққұнды ақуыздың, полиқанықпаған май қышқылдарының және басқа да
биологиялық белсенді компоненттердің маңызды көзі ретінде қарастырылады
{[}1,2{]}.

Балық өнімдерінің құрылымында ready-to-eat (RTE) санатына жататын
өнімдер ерекше орын алады. Бұл өнімдер қосымша термиялық өңдеусіз
тұтынылатындықтан, олардың микробиологиялық қауіпсіздігін қамтамасыз
етуге, өндірістің санитарлық жағдайларын бақылауға, сақтау мен
тасымалдау кезіндегі температуралық режимдерді сақтауға, сондай-ақ
шикізат пен дайын өнімнің қадағалануын қамтамасыз етуге жоғары талаптар
қойылады {[}3,4{]}.

Ферменттелген балық өнімдері күрделі микробиологиялық экожүйемен
сипатталатын тағам өнімдерінің жеке санатын құрайды. Стерилизацияланған
немесе термиялық өңдеуден өткен өнімдерден айырмашылығы, ферменттелген
өнімдердің қауіпсіздігі тек өндірістің технологиялық параметрлерімен
ғана емес, сонымен қатар микроорганизмдердің дамуын бақылауды, суық
тізбектің сақталуын және шикізаттың шығу тегін бақылауды қамтитын
нормативтік бақылау жүйесінің тиімділігімен анықталады {[}5,6{]}.

Тилапия (Tilapia nilotica) әлемдік аквакультураның ең қарқынды дамып
келе жатқан объектілерінің бірі болып табылады. Оның етінің жоғары
биологиялық құндылығы, қолайлы аминқышқылдық құрамы және технологиялық
әмбебаптығы бұл балықты әртүрлі өңделген балық өнімдерін өндіру үшін
перспективалы шикізат ретінде қарастыруға мүмкіндік береді. Аквакультура
өнімдерінің халықаралық саудасы артып отырған жағдайда тилапиядан
алынатын өнімдердің халықаралық азық-түлік қауіпсіздігі талаптарына және
әртүрлі елдердің нормативтік стандарттарына сәйкестігін қамтамасыз ету
ерекше маңызға ие болуда {[}1-6{]}.

Еуропалық одақтағы азық-түлік қауіпсіздігін реттеу жүйесі тәуекелге
негізделген тәсілге және «фермадан тұтынушыға дейін» қағидатына
негізделген. Бұл жүйе өндірістің гигиенасына қойылатын талаптарды,
микробиологиялық қауіпсіздік критерийлерін, ресми бақылауды, өнімнің
қадағалануын және сақтау мерзімін ғылыми негіздеуді қамтитын өнімнің
бүкіл өндірістік-тасымалдау тізбегі бойынша кешенді бақылауды
қарастырады. Бұл модель халықаралық деңгейде азық-түлік қауіпсіздігін
реттеудің ең дамыған және тиімді жүйелерінің бірі болып саналады {[}7 -
10{]}.

Еуразиялық экономикалық одақ елдерінде азық-түлік өнімдерінің
қауіпсіздігін реттеу тұтынушылардың денсаулығын қорғауға және ортақ
нарықта өнім айналымының бірыңғай ережелерін қалыптастыруға бағытталған
техникалық регламенттер негізінде жүзеге асырылады. Негізгі нормативтік
құжаттарға КО ТР 021/2011 «Тағам өнімдерінің қауіпсіздігі туралы», КО ТР
022/2011 «Тағам өнімдерінің таңбалануы туралы» және КО ТР 040/2016
«Балық және балық өнімдерінің қауіпсіздігі туралы» техникалық
регламенттері жатады {[}10{]}.

Алайда халықаралық сауданың дамуы және санитарлық-фитосанитарлық
талаптардың күшеюі жағдайында ұлттық нормативтік талаптарды халықаралық
стандарттармен және Codex Alimentarius ұсынымдарымен одан әрі үйлестіру
қажеттілігі туындап отыр {[}3{]}.

Заманауи зерттеулер әртүрлі елдердің нормативтік талаптарының формалды
түрде ұқсас болуы олардың іс жүзіндегі функционалдық сәйкестігін әрқашан
қамтамасыз ете бермейтінін көрсетеді. Айтарлықтай айырмашылықтар бақылау
жүйелерінің цифрландыру деңгейінде, зертханалық мониторингтің
қарқындылығында, антифрод-бақылау механизмдерінің тиімділігінде және
реттеуші органдардың азық-түлік қауіпсіздігі талаптарының бұзылуына
жедел әрекет ету қабілетінде байқалуы мүмкін {[}5-7{]}.

Осыған байланысты тағам қауіпсіздігін қамтамасыз етудің нормативтік
жүйелерінің үйлесімділік деңгейін бағалаудың сандық құралдарын әзірлеу
өзекті ғылыми міндеттердің бірі болып табылады. Мұндай құралдардың бірі
ретінде нормативтік реттеудің әртүрлі элементтерін салыстыруға және
ұлттық және халықаралық талаптар арасындағы негізгі реттеуші
айырмашылықтарды анықтауға мүмкіндік беретін интегралды үйлесімділік
индексі қолданылуы мүмкін {[}7-10{]}.

Осы зерттеудің мақсаты - Еуропалық одақ талаптарымен және Codex
Alimentarius халықаралық стандарттарымен салыстыру негізінде тилапиядан
(Tilapia nilotica) алынатын өнімдерді қоса алғанда, балық өнімдерінің
қауіпсіздігін нормативтік реттеудің үйлесімділік деңгейін бағалаудың
интегралды моделін әзірлеу, негізгі реттеуші блоктар бойынша gap-талдау
жүргізу және балық өнімдерінің қауіпсіздігін нормативтік қамтамасыз ету
жүйесін жетілдіру бойынша ғылыми негізделген ұсынымдар қалыптастыру.

{\bfseries Материалдар мен әдістер.} Зерттеу нысаны ретінде Қазақстан
Республикасындағы және Еуразиялық экономикалық одақ (ЕАЭО) елдеріндегі
балық өнімдерінің сапасы мен қауіпсіздігін нормативтік реттеу жүйесі,
сондай-ақ оның азық-түлік қауіпсіздігін қамтамасыз ету саласындағы
халықаралық талаптарға сәйкестік деңгейі қарастырылды. Ерекше назар
балық өнімдерінің өндірістік-өткізу тізбегінің негізгі кезеңдерін
реттейтін нормативтік механизмдерді талдауға аударылды. Бұл кезеңдерге
шикізат өндірісі, өңдеу, сақтау, тасымалдау және дайын өнімді өткізу
процестері кіреді.

Зерттеудің нормативтік-құқықтық негізін балық өнімдерінің қауіпсіздігін
реттейтін халықаралық, өңірлік және ұлттық құжаттар құрады. Негізгі
дереккөздер ретінде Codex Alimentarius стандарттары, Еуропалық одақтың
нормативтік актілері, Еуразиялық экономикалық одақтың техникалық
регламенттері, сондай-ақ Қазақстан Республикасының қолданыстағы
нормативтік-құқықтық құжаттары пайдаланылды.

Атап айтқанда, зерттеу барысында келесі нормативтік құжаттарға талдау
жүргізілді:

- ЕАЭО техникалық регламенті КО ТР 021/2011 «Тағам өнімдерінің
қауіпсіздігі туралы»;

- ЕАЭО техникалық регламенті КО ТР 040/2016 «Балық және балық
өнімдерінің қауіпсіздігі туралы»;

- ЕАЭО техникалық регламенті КО ТР 022/2011 «Тағам өнімдерінің
таңбалануы туралы»;

- балық және балық өнімдерін өңдеу қауіпсіздігін реттейтін Codex
Alimentarius халықаралық стандарттары;

- азық-түлік өнімдерінің гигиенасы, микробиологиялық қауіпсіздік
критерийлері және ресми бақылау саласындағы Еуропалық одақ
регламенттері.

Зерттеудің әдіснамалық негізін балық өнімдерінің қауіпсіздігін реттеудің
қолданыстағы нормативтік жүйесін жүйелі бағалауға мүмкіндік беретін
өзара толықтыратын ғылыми әдістер кешені құрады. Зерттеу барысында
салыстырмалы-құқықтық талдау, нормативтік талаптарды жүйелік талдау,
нормативтік құжаттардың контент-талдауы, сондай-ақ сараптамалық бағалау
және нәтижелерді сандық интерпретациялау әдістері қолданылды.

Зерттеудің бірінші кезеңінде балық өнімдерінің қауіпсіздігін реттейтін
нормативтік құжаттарға контент-талдау жүргізілді. Бұл кезең балық
өнімдерін өндіру, сақтау және өткізу кезінде қойылатын негізгі
талаптарды жүйелеуге мүмкіндік берді. Атап айтқанда, микробиологиялық
қауіпсіздік критерийлері, тағам өнеркәсібі кәсіпорындарына қойылатын
санитарлық-гигиеналық талаптар, суық тізбекті сақтау шарттары және
мемлекеттік бақылау рәсімдері қарастырылды.

Екінші кезеңде халықаралық стандарттар мен ЕАЭО-ның қолданыстағы
нормативтік жүйесі арасындағы ұқсастықтар мен айырмашылықтарды анықтауға
бағытталған салыстырмалы талдау жүргізілді. Салыстыру балық өнімдерінің
қауіпсіздігін реттеудің келесі негізгі бағыттары бойынша жүзеге
асырылды:

- өнімнің микробиологиялық қауіпсіздігі;

- өндіріс және өңдеу гигиенасына қойылатын талаптар;

- сақтау және тасымалдау кезінде температуралық режимдерді және суық
тізбекті бақылау;

- өнімнің қадағалануын қамтамасыз ету және шикізаттың шығу тегін
сәйкестендіру;

- таңбалау және тұтынушыларды ақпараттандыру талаптары;

- мемлекеттік бақылау және инспекциялық қадағалау рәсімдері.

Ұлттық талаптардың халықаралық стандарттарға сәйкестік деңгейін сандық
бағалау үшін нормативтік талаптардың үйлесімділігін интегралды бағалау
әдісі қолданылды. Осы тәсіл шеңберінде ұлттық реттеу жүйесінің
халықаралық стандарттар талаптарына сәйкестік деңгейін бағалауға
мүмкіндік беретін IH-FISH (Index of Harmonization for Fish Safety
Regulation) интегралды үйлесімділік индексінің моделі әзірленді.

Бағалау бірнеше негізгі реттеуші блоктар бойынша жүргізілді, олардың
әрқайсысына халықаралық талаптарға сәйкестік деңгейіне байланысты балдық
баға берілді.

Бағалау үшін үш деңгейлі шкала қолданылды:

0 - нормативтік жүйеде ұқсас талаптардың болмауы;

1 - халықаралық талаптарға ішінара сәйкестік;

2 - халықаралық стандарттарға толық сәйкестік.

Әрбір реттеуші блоктың салыстырмалы маңыздылығын анықтау үшін тағам
қауіпсіздігі, стандарттау және нормативтік реттеу саласындағы мамандар
қатысқан сараптамалық ранжирлеу әдісі қолданылды. Сараптамалық бағалау
балық өнімдерінің қауіпсіздігін реттеу жүйесінің әртүрлі элементтерінің
салмақтық коэффициенттерін анықтауға мүмкіндік берді.

Қосымша ретінде ЕАЭО-ның қолданыстағы реттеу жүйесі мен халықаралық
стандарттар арасындағы нормативтік айырмашылық деңгейін сандық тұрғыдан
анықтауға мүмкіндік беретін gap-талдау әдісі қолданылды. Бұл әдіс
нормативтік реттеудің одан әрі үйлестіруді және жетілдіруді талап ететін
ең проблемалық бағыттарын анықтауға мүмкіндік берді.

Алынған деректерді өңдеу және түсіндіру үшін зерттеу нәтижелерін
жүйелеуге және балық өнімдерінің қауіпсіздігін нормативтік реттеу
жүйесінің дамуындағы негізгі үрдістерді анықтауға мүмкіндік беретін
сипаттамалық статистика әдістері қолданылды.

Кешенді әдіснамалық тәсілді қолдану балық өнімдерінің қауіпсіздігін
реттеудің қолданыстағы жүйесін жан-жақты талдауға және оны азық-түлік
қауіпсіздігінің халықаралық стандарттарымен үйлестіру жағдайында одан
әрі жетілдірудің негізгі бағыттарын анықтауға мүмкіндік берді.

{\bfseries Нәтижелер және талқылау.} Жүргізілген зерттеу балық өнімдерінің
сапасы мен қауіпсіздігін нормативтік реттеудің қолданыстағы жүйесіне
кешенді талдау жүргізуге және оның халықаралық стандарттармен
үйлесімділік деңгейін анықтауға мүмкіндік берді. Зерттеу барысында
Еуропалық одақтың нормативтік талаптары, Codex Alimentarius халықаралық
стандарттары және Еуразиялық экономикалық одақтың техникалық
регламенттері салыстырмалы түрде қарастырылды.

Талдау нәтижелері көрсеткендей, қарастырылған реттеу жүйелерінің барлығы
азық-түлік қауіпсіздігін қамтамасыз етудің ортақ негізгі қағидаттарына
негізделген. Мұндай қағидаттарға микробиологиялық қауіптердің алдын алу,
өндірістің санитарлық-гигиеналық жағдайларын қамтамасыз ету, сапаны
бақылау жүйелерін енгізу және тағам өнеркәсібі кәсіпорындарында
гигиеналық талаптарды сақтау жатады. Сонымен қатар нормативтік
талаптардың нақтылану деңгейі мен оларды практикалық іске асыру
механизмдері арасында айтарлықтай айырмашылықтар байқалады.

Балық өнімдерінің қауіпсіздігін нормативтік реттеудің ең дамыған жүйесі
Еуропалық одақ елдерінде қалыптасқан. Еуропалық модель «фермадан
тұтынушыға дейін» қағидатына негізделген және өнімнің өндірістік-өткізу
тізбегінің барлық кезеңдерінде қауіпсіздігін кешенді бақылауды
қарастырады. Бұл жүйеде өнімнің қадағалануына, сақтау мерзімінің ғылыми
негізделуіне, цифрлық мониторинг жүйелеріне және тәуекелге негізделген
мемлекеттік бақылауға ерекше көңіл бөлінеді.

Codex Alimentarius халықаралық стандарттары азық-түлік қауіпсіздігін
қамтамасыз ету бойынша жаһандық ұсынымдар жүйесін қалыптастырады. Бұл
стандарттардың негізгі мақсаты -- әртүрлі елдердің талаптарын үйлестіру
және халықаралық саудада тағам өнімдерінің қауіпсіздігін қамтамасыз ету.
Ұсынымдық сипатта болғанымен, Codex стандарттары ұлттық нормативтік
құжаттарды әзірлеу кезінде кеңінен қолданылады және азық-түлік
қауіпсіздігін реттеудің заманауи жүйесін қалыптастыруда бағдар ретінде
қызмет етеді.

Еуразиялық экономикалық одақ елдерінде балық өнімдерінің қауіпсіздігін
реттеу өнімді өндіру, сақтау, тасымалдау және өткізуге қойылатын
міндетті талаптарды белгілейтін техникалық регламенттер негізінде жүзеге
асырылады. Негізгі нормативтік құжаттар ретінде КО ТР 021/2011 «Тағам
өнімдерінің қауіпсіздігі туралы» және КО ТР 040/2016 «Балық және балық
өнімдерінің қауіпсіздігі туралы» техникалық регламенттері қарастырылады.
Бұл құжаттар балық өнімдерінің сапасы мен қауіпсіздігіне қойылатын
негізгі талаптарды айқындайды.

Нормативтік талаптарды салыстырмалы талдау нәтижелері 1-кестеде
көрсетілген. Кестеден көрініп тұрғандай, нормативтік жүйелер арасындағы
ең жоғары үйлесімділік деңгейі өнімнің микробиологиялық қауіпсіздігі
және өндірістің санитарлық талаптары саласында байқалады. ЕАЭО
техникалық регламенттерінде патогенді микроорганизмдердің рұқсат етілген
деңгейлері белгіленген және олар халықаралық стандарттар мен Еуропалық
одақтың нормативтік талаптарымен салыстырмалы деңгейде.

Сонымен қатар өнімнің қадағалануын қамтамасыз ету, суық тізбекті бақылау
және сақтау мерзімін ғылыми негіздеу салаларында белгілі бір
айырмашылықтар анықталды. Еуропалық одақ елдерінде бұл талаптар
анағұрлым егжей-тегжейлі регламенттелген және өнім қозғалысын бақылаудың
цифрлық жүйелері мен заманауи мониторинг құралдарын енгізумен қатар
жүреді.

Балық өнімдерінің әлемдік нарығында ерекше мәселе ретінде өнімді
фальсификациялау жағдайлары қарастырылады. Бұған балық түрлерін
ауыстыру, өнімнің шыққан елін дұрыс көрсетпеу және сақтау шарттарын бұзу
сияқты заңбұзушылықтар жатады. Мұндай бұзушылықтардың алдын алу
мақсатында Еуропалық одақ елдерінде шикізатты сәйкестендірудің заманауи
әдістері, соның ішінде молекулалық-генетикалық талдау әдістері және
өнімнің қадағалануын қамтамасыз ететін цифрлық жүйелер кеңінен
қолданылады. Ал ЕАЭО-ның нормативтік жүйесінде мұндай механизмдер
қазіргі уақытта даму кезеңінде. Нормативтік талаптарды салыстырмалы
талдау нәтижелері 1-кестеде ұсынылған.

Кестеден көрініп тұрғандай, үйлесімділіктің ең жоғары деңгейі өнімнің
микробиологиялық қауіпсіздігі және өндірістің санитарлық талаптары
саласында байқалады. Орташа сәйкестік деңгейі суық тізбекті бақылау,
өнімнің қадағалануын қамтамасыз ету және таңбалау талаптары саласында
анықталды. Үйлесімділіктің ең төмен деңгейі антифрод-бақылау саласында
байқалады, бұл балық өнімдерін фальсификациялаудың алдын алу
механизмдерін одан әрі жетілдіру қажеттілігін көрсетеді. Реттеу жүйелері
арасындағы нормативтік айырмашылықтарды неғұрлым нақты анықтау үшін
gap-талдау жүргізілді, оның нәтижелері 2-кестеде ұсынылған.
\end{multicols}
\vspace{-2em}
\tcap{1-кесте. Балық өнімдерінің қауіпсіздігіне қойылатын нормативтік
талаптардың үйлесімділігін салыстырмалы талдау}
\begin{longtblr}[
  label = none,
  entry = none,
]{
  width = \linewidth,
  colspec = {Q[110]Q[325]Q[206]Q[208]Q[87]},
  cells = {font = \small},
  row{1} = {c,font=\bfseries\small},
  hlines,
  vlines,
}
Реттеуші блок & Еуропалық одақ & Codex Alimentarius & ЕАЭО & Үйлесім\-ділік деңгейі\\
Микробио\-логиялық қауіпсіздік & Микробиологиялық қауіпсіздік критерийлері толық регламенттелген және тұрақты зертханалық бақылау жүргізіледі & Микробиологиялық критерийлер бойынша халықаралық ұсынымдар & Патогенді микроорганизмдердің рұқсат етілген деңгейлері белгіленген & Жоғары\\
Өндірістік гигиена & HACCP жүйесін міндетті енгізу және қатаң мемлекеттік бақылау & Тағам өнімдерінің гигиенасының жалпы қағидаттары & HACCP жүйесі кәсіпорындарда енгізілуде & Жоғары\\
Суық тізбекті бақылау & Температуралық режимдерді қатаң мониторингтеу және логистиканы цифрлық бақылау & Сақтау және тасымалдау температурасы бойынша ұсынымдар & Ішінара регламенттелген & Орташа\\
Өнімнің қадағалану жүйесі & Жеткізу тізбегінің барлық кезеңдерінде толық traceability жүйесі & Ұсынымдық талаптар & Қадағалану жүйесі ішінара енгізілген & Орташа\\
Өнімді таңбалау & Өнімді таңбалау және тұтынушыларды ақпараттандыру бойынша толық талаптар & Таңбалау бойынша жалпы ұсынымдар & Таңбалауға қойылатын негізгі талаптар белгіленген & Орташа\\
Антифрод-бақылау & Өнім түрін және шыққан елін алмастыруды бақылаудың дамыған жүйесі & Шектеулі ұсынымдар & Іс жүзінде жоқ & Төмен\\
Сақтау мерзімін валидациялау & Сақтау мерзімін ғылыми негізде міндетті түрде растау (shelf-life validation) & Сақтау мерзімдерін белгілеу бойынша ұсынымдар & Ішінара реттелген & Орташа
\end{longtblr}
\vspace{-2em}
\tcap{2-кесте. Нормативтік айырмашылықтардың gap-талдауы}
\begin{longtblr}[
  label = none,
  entry = none,
]{
  cells = {c},
  cells = {font = \small},
  row{1} = {font=\bfseries\small},
  hlines,
  vlines,
}
Реттеуші блок & Салмақ коэффициенті & Сәйкестік деңгейі & Қорытынды\\
Микробиологиялық қауіпсіздік & 0,25 & 2 & 0,50\\
Суық тізбек & 0,15 & 1 & 0,15\\
Өнімнің қадағалану жүйесі & 0,20 & 1 & 0,20\\
Антифрод-бақылау & 0,10 & 0 & 0\\
Мемлекеттік бақылау & 0,15 & 2 & 0,30\\
Сақтау мерзімін валидациялау & 0,15 & 1 & 0,15
\end{longtblr}
\vspace{-2em}
\begin{multicols}{2}
2-кесте деректері бойынша ең үлкен нормативтік айырмашылық
антифрод-бақылау саласында байқалады, мұнда gap көрсеткішінің мәні 2-ге
тең. Бұл Еуропалық одақ талаптары мен ЕАЭО-ның қолданыстағы нормативтік
жүйесі арасында айтарлықтай айырмашылық бар екенін көрсетеді.
Айтарлықтай айырмашылықтар сондай-ақ өнімнің қадағалануын қамтамасыз
ету, суық тізбекті бақылау және сақтау мерзімін ғылыми негіздеу
салаларында анықталды.

Нормативтік талаптардың үйлесімділік деңгейін сандық тұрғыдан бағалау
үшін үйлесімділіктің интегралды индексі есептелді, оның нәтижелері
3-кестеде ұсынылған.
\end{multicols}

\tcap{3-кесте. Нормативтік жүйенің үйлесімділігін интегралды бағалау}
\begin{longtblr}[
  label = none,
  entry = none,
]{
  width = \linewidth,
  colspec = {Q[327]Q[258]Q[204]Q[148]},
  cells = {c},
  cells = {font = \small},
  row{1} = {font=\bfseries\small},
  hlines,
  vlines,
}
Реттеуші блок & Салмақ коэффициенті & Сәйкестік деңгейі & Қорытынды\\
Микробиологиялық қауіпсіздік & 0,25 & 2 & 0,50\\
Суық тізбек & 0,15 & 1 & 0,15\\
Қадағалану (traceability) & 0,20 & 1 & 0,20\\
Антифрод-бақылау & 0,10 & 0 & 0\\
Мемлекеттік бақылау & 0,15 & 2 & 0,30\\
Сақтау мерзімін валидациялау & 0,15 & 1 & 0,15
\end{longtblr}

\begin{multicols}{2}
Жоғарыда үйлесімділіктің интегралды бағасына ең үлкен үлесті
микробиологиялық қауіпсіздік және мемлекеттік бақылау көрсеткіштері
қосады. Сонымен қатар антифрод-бақылау және өнімнің қадағалануын
қамтамасыз ету салаларында салыстырмалы түрде төмен мәндер байқалады.

Үйлесімділіктің интегралды индексін есептеу нәтижелері ЕАЭО нормативтік
жүйесінің халықаралық талаптарға сәйкестік деңгейін орташа деңгейде
бағалауға болатынын көрсетті. Алынған индекс мәні қолданыстағы
нормативтік жүйе жалпы алғанда халықаралық стандарттарға сәйкес
келетінін, алайда реттеудің жекелеген элементтерін одан әрі жетілдіру
қажеттігін көрсетеді.

Нормативтік талаптардың үйлесімділік деңгейінің графикалық көрінісі
1-суретте ұсынылған.
\end{multicols}

\fig[0.7\textwidth]{p2/image11}[1-сурет. Нормативтік талаптардың
үйлесімділік деңгейінің гистограммасы]

\begin{multicols}{2}
Гистограммада микробиологиялық қауіпсіздік және өндірістің санитарлық
талаптары саласында үйлесімділіктің ең жоғары деңгейі байқалатынын айқын
көрсетеді. Сонымен қатар ең үлкен айырмашылықтар антифрод-бақылау
саласында анықталған, бұл жүргізілген gap-талдау нәтижелерін растайды.

Балық өнімдерінің қауіпсіздігін қамтамасыз етудің кешенді тәсілі
2-суретте ұсынылған, онда балық өнімдерінің қауіпсіздігін реттеу
жүйесінің концептуалдық схемасы көрсетілген.

\fig[\columnwidth]{p2/image12}[2-сурет. Балық өнімдерінің
қауіпсіздігін реттеу жүйесінің концептуалдық схемасы]
\vspace{-1em}
Суретте балық өнімдерінің қауіпсіздігін қамтамасыз етудің заманауи
моделі өндірістік-өткізу тізбегінің барлық кезеңдерін - шикізатты
өндіруден бастап соңғы тұтынушыға дейін -- дәйекті бақылауға
негізделген. Бастапқы кезеңде балықты аулау немесе оны аквакультура
жағдайында өсіру жүзеге асырылады. Бұл кезең өндірістің экологиялық
жағдайларын, су сапасын, жемшөп базасын және балық шикізатының
ветеринариялық жағдайын бақылауды қамтиды.

Келесі кезең -- шикізатты өңдеу. Бұл кезең балық шикізатын дайындау, оны
ұсақтау, жартылай фабрикаттарды қалыптастыру және дайын өнім өндіру
сияқты технологиялық операцияларды қамтиды. Осы кезеңде
санитарлық-гигиеналық талаптарды сақтау және өндірістің технологиялық
параметрлерін орындау ерекше маңызға ие.

Тағам өнімдерінің қауіпсіздігін қамтамасыз ету жүйесінің негізгі
элементі HACCP (Hazard Analysis and Critical Control Points) жүйесін
енгізу болып табылады. Бұл жүйе өндірістің барлық кезеңдерінде ықтимал
биологиялық, химиялық және физикалық қауіптерді анықтауға, бағалауға
және басқаруға бағытталған.

Бақылау жүйесінің маңызды құрамдас бөлігі -- суық тізбектің
үздіксіздігін қамтамасыз ету. Бұл балық өнімдерін сақтау және тасымалдау
кезінде оңтайлы температуралық режимді сақтауға мүмкіндік береді.
Температуралық талаптарды сақтау патогенді микроорганизмдердің дамуын
болдырмаудың және өнім сапасын сақтаудың негізгі факторы болып табылады.

{\bfseries Қорытынды.} Жүргізілген зерттеу балық өнімдерінің сапасы мен
қауіпсіздігін нормативтік реттеу жүйесіне кешенді бағалау жүргізуге және
оның азық-түлік қауіпсіздігі саласындағы халықаралық стандарттармен
үйлесімділік деңгейін анықтауға мүмкіндік берді. Зерттеу барысында
Еуропалық одақтың талаптары, Codex Alimentarius халықаралық стандарттары
және Еуразиялық экономикалық одақтың техникалық регламенттері
салыстырылды.

Зерттеу нәтижелері қарастырылған нормативтік жүйелердің азық-түлік
қауіпсіздігін қамтамасыз етудің ортақ қағидаттарына негізделетінін
көрсетті. Оларға микробиологиялық қауіптерді бақылау, өндірістің
санитарлық-гигиеналық талаптарын сақтау және HACCP қағидаттарына
негізделген азық-түлік қауіпсіздігін басқару жүйесін енгізу жатады.
Сонымен бірге нормативтік талаптардың нақтылану деңгейі мен оларды
практикалық іске асыру механизмдері арасында елеулі айырмашылықтар
байқалады.

Талдау нәтижелері нормативтік талаптардың ең жоғары үйлесімділігі
өнімнің микробиологиялық қауіпсіздігі және өндірістің санитарлық
талаптары саласында байқалатынын көрсетті. Сонымен қатар суық тізбекті
бақылау, өнімнің қадағалану жүйелері, сақтау мерзімін ғылыми негіздеу
және балық өнімдерінің фальсификациясының алдын алу механизмдерін реттеу
салаларында айырмашылықтар анықталды.

Жүргізілген gap-талдау антифрод-бақылау және өнімнің цифрлық қадағалану
жүйелері саласында ең үлкен нормативтік айырмашылықтар бар екенін
көрсетті. IH-FISH (Index of Harmonization for Fish Safety Regulation)
индексіне негізделген үйлесімділікті бағалаудың интегралды моделі ЕАЭО
нормативтік жүйесінің халықаралық талаптарға сәйкестік деңгейін сандық
тұрғыдан анықтауға мүмкіндік берді және ол орташа деңгейде екені
анықталды.

Алынған нәтижелер суық тізбекті цифрлық мониторингтеу жүйелерін дамыту,
өнімнің қадағалану жүйелерін жетілдіру және балық өнімдерінің
фальсификациясына қарсы шараларды күшейту қажеттігін көрсетеді.
Ұсынылған үйлесімділікті бағалау моделі азық-түлік қауіпсіздігін
реттеудің ұлттық жүйелерінің тиімділігін талдау үшін, сондай-ақ балық
өнімдерінің халықаралық нарықтағы экспорттық бәсекеге қабілеттілігін
арттыру мақсатында қолданылуы мүмкін.
\end{multicols}

\begin{center}
{\bfseries Әдебиеттер}
\end{center}

\begin{refs}
1. FAO. The State of World Fisheries and Aquaculture 2022: Towards Blue
Transformation Rome: FAO - 2022.-266 p. ISBN 978-92-5-136364-5. DOI
\href{https://doi.org/10.4060/cc0461en}{10.4060/cc0461en}.

2. FAO. The State of World Fisheries and Aquaculture 2024: Blue
Transformation in Action. Rome: FAO, 2024. - 264 p. ISBN
978-92-5-138763-4. DOI 10.4060/cd0683en.

3. FAO, WHO. Code of Practice for Fish and Fishery Products (CAC/RCP
52--2003) - Codex Alimentarius Commission. -- Rome: FAO/WHO.
--372 p. ISBN 978-92-4-001318-6

4. Yin M., Wang X. New strategies to improve the quality and safety of
seafoods and the efficient utilization of their by-products // Foods.
-2026.-Vol.15(2):190. DOI
\href{https://doi.org/10.3390/foods15020190}{10.3390/foods15020190}.

5. Issa-Zacharia A., Msangi R., Mboya G. Chemical and microbiological
quality of fish products: regulatory aspects // International Journal of
Microbiology and Biotechnology. -2025.-Vol10(3). -P.73-90. - DOI
\href{https://doi.org/10.11648/j.ijmb.20251003.12}{10.11648/j.ijmb.20251003.12}.

6. Visciano P. Environmental contaminants in fish products: food safety
issues and remediation strategies // Foods. - 2024. -Vol.13(21):3511.
DOI \href{https://doi.org/10.3390/foods13213511}{10.3390/foods13213511}.

7. Ryder J., Karunasagar I., Ababouch L. (eds.). Assessment and
Management of Seafood Safety and Quality: Current Practices and Emerging
Issues -- FAO Fisheries and Aquaculture Technical Rome: FAO, 2014. - 455
p. ISBN 978-92-5-107511-1

8. Eegulation (EC) No.178/2002 of the European Parliament and of the
Council laying down the general principles and requirements of food law,
establishing the European Food Safety Authority and laying down
procedures in matters of food safety/ Official Journal L 31, 1 February
2002, pp.1-24.- Source:
FAO, \href{https://www.fao.org/faolex/results/details/en/c/LEX-FAOC034771}{FAOLEX}.
\href{https://www.ecolex.org/details/legislation/regulation-ec-no-1782002-of-the-european-parliament-and-of-the-council-laying-down-the-general-principles-and-requirements-of-food-law-establishing-the-european-food-safety-authority-and-laying-down-procedures-in-matters-of-food-safety-lex-faoc034771/?type=legislation&xsubjects=Energy&page=170}{URL: https://www.ecolex.org} .- Date of address:-03.01.2026

9. Regulation (EC) No 852/2004 of the European parliament and of the
council of 29 April 2004 on the hygiene of
foodstuffs \url{https://eur-lex.europa.eu/eli/reg/2004/852/oj/eng} -
Date of address:-03.01.2026.

10. O tehnicheskom reglamente Evrazijskogo jekonomicheskogo sojuza «O
bezopasnosti ryby i rybnoj produkcii». Reshenie Soveta Evrazijskoj
jekonomicheskoj komissii ot 18.10.2016 g. №162.
\href{https://adilet.zan.kz/rus/docs/H16EV000162\%20-}{https://adilet.zan.kz/rus/docs/H16EV000162
-} Data obrashhenija:03.01.2026.
\end{refs}

\begin{info}
\hspace{1em}\emph{{\bfseries Авторлар туралы мәліметтер}}

Курмангалиева Д.Б. - техника ғылымдарының докторы, профессор м.а., С.
Сейфуллин атындағы Қазақ агротехникалық зерттеу университеті, Астана,
Қазақстан, e-mail: d.kurmangaliyeva@kazatu.edu.kz;

Дарибаева А.А. - PhD докторанты, С. Сейфуллин атындағы Қазақ
агротехникалық зерттеу университеті, Астана., Қазақстан, e-mail:
aigul\_daribaeva@mail.ru; Юсупова Юсупова Г.Т. - PhD, аға оқытушы,
«Көлік техника, технология және стандарттау, сертификаттау және
метрология» БББТ, С.Сейфуллин атындағы Қазақ агротехникалық зерттеу
университеті, Астана қ., Қазақстан, е-mail: gauhar\_20\_90@mail.ru;

Шадьярова Ж.К.- PhD, аға оқытушы, С.Сейфуллин атындағы Қазақ
агротехникалық зерттеу университеті, Астана, Қазақстан, е-mail:
zhazirashadyarova@gmail.com;

Жуматова Г.Г. - экология магистрі, Халел Досмұхамедова атындағы Атырау
университеті. Атырау, Қазақстан, e-mail: gulshat\_zhumatova@mail.ru;

\hspace{1em}\emph{{\bfseries Information about the authors}}

Kurmangaliyeva D.- Doctor of Technical Sciences, Acting Professor, S.
Seifullin Kazakh Agrotechnical Research University, Astana,
Kazakhstan; e-mail: d.kurmangaliyeva@kazatu.edu.kz

Daribayeva А. - PhD Doctoral Student, S. Seifullin Kazakh
Agrotechnical Research University, Astana, Kazakhstan; e-mail:
aigul\_daribaeva@mail.ru;

Yussupova G. - PhD, Senior Lecturer, S.Seifullin Kazakh Agrotechnical
Research University, Astana, Kazakhstan; e-mail:
gauhar\_20\_90@mail.ru;

Shadyarova Zh. - PhD, Senior Lecturer, S.Seifullin Kazakh
Agrotechnical Research University, Astana, Kazakhstan, е-mail:
zhazira\_shadyarova@bk.ru

Zhumatova G.-- Master of Ecology, H. Dosmukhamedova Atyrau University,
Atyrau, Kazakhstan, e-mail: gulshat\_zhumatova@mail.ru.
\end{info}
